\documentclass{article}
\usepackage[margin=1in]{geometry}
\usepackage{amsmath,amssymb,amsthm}
\usepackage{hyperref}
\usepackage{tcolorbox}
\tcbuselibrary{breakable}

\newtcolorbox{accept}[1][]{colback=green!5,colframe=green!60!black,
  fonttitle=\bfseries,title={Accepted: #1},breakable}
\newtcolorbox{partial}[1][]{colback=yellow!5,colframe=orange!60!black,
  fonttitle=\bfseries,title={Partially Accepted: #1},breakable}
\newtcolorbox{correct}[1][]{colback=blue!5,colframe=blue!60!black,
  fonttitle=\bfseries,title={Key Insight / Correction: #1},breakable}
\newtcolorbox{dispute}[1][]{colback=red!3,colframe=red!50!black,
  fonttitle=\bfseries,title={Contested / Nuanced: #1},breakable}

\title{Response to Reviewer Op46MX}
\date{April 2026}
\begin{document}
\maketitle

\section{Overall Assessment}

Op46MX provides an extremely detailed review with several original mathematical
contributions. The most important is the \textbf{corrected curvature proxy
scaling law} which resolves the apparent inconsistency between Equation~(16)
and Remark~8. The review also correctly identifies the torus data bug, the edge
KL term, and sphere sampling non-uniformity. We accept all critical findings.

\section{Critical Weaknesses -- Response}

\subsection{W1: Assumption A4' (Critical)}

\begin{accept}[Accepted -- Sobolev interpolation lemma (Lemma 1 of Op46MX)]
Op46MX provides the clearest statement and proof of the capacity condition:
\textbf{Lemma 1 (Jacobian control via Sobolev interpolation)}: under bounded
Hessian and reconstruction accuracy $\delta$ at training points,
\[
  \|J_{f^*}(z) - J_{f_{\mathrm{true}}}(z)\|_{\mathrm{op}} \leq \frac{2\delta}{r_n} + \kappa r_n
\]
and \textbf{Corollary 1}: if $\delta_{\mathrm{rec}} = O(r_n^2)$ then A4' holds
with $\|G^* - g\|_{\mathrm{op}} = O(r_n)$.

\textbf{Evaluation of the proof}: The proof is correct. The key step uses the
Taylor expansion and reconstruction accuracy at adjacent training points to bound
the directional Jacobian difference. The proof is cleaner than the finite-
difference argument we had. We will incorporate this lemma and corollary into
the paper, with attribution.

\textbf{Concerning the decoder capacity requirement}: Op46MX notes that for $d=2$,
$n=3000$, the Barron-type bound requires hidden width $\geq 400$, and our
hidden widths 128/256 ``may be marginal.'' This is a fair empirical concern.
We will add an empirical verification (plotting $\|G^* - g\|$ vs.\ $r_n$
for the sphere where the true metric is known analytically) to validate
that A4' holds approximately in our experimental regime.
\end{accept}

\subsection{W2: Curvature Proxy Prediction Inconsistency (Critical)}

\begin{correct}[The most important theoretical insight of this review]
Op46MX provides the derivation showing that:
\[
  \hat{\kappa} = \frac{\|c_{ijk}\|}{\|\ell_{ij}\|\|\ell_{ik}\|}
  \approx \frac{2}{R} \cdot \frac{1}{\sqrt{G/3}} = \frac{2}{R}\sqrt{\frac{3}{G}}
\]
for an uncalibrated random embedding with $A \sim \mathcal{N}(0, 1/\sqrt{3})$.
For $R=1$, $G=50$: $\hat{\kappa}^* = 0.4899$.

This directly contradicts Eq.~(16)'s claim of $\hat{\kappa} \to 2/R = 2.000$.

\textbf{Mathematical evaluation}: Op46MX's derivation is correct. In the
ambient space, both the numerator $\|c_{ijk}\|$ and denominator
$\|\ell_{ij}\|\|\ell_{ik}\|$ are scaled by factors of $\sqrt{G/3}$:
\[
  \hat{\kappa} = \frac{\sqrt{G/3} \cdot 2\varepsilon^2/R}{(G/3)\varepsilon^2} = \frac{2}{R\sqrt{G/3}}
\]
The numerator scales as $\sqrt{G/3}$ while the denominator scales as $G/3$,
leaving a $1/\sqrt{G/3}$ factor that Eq.~(16) omits.

\textbf{However} -- and this is the KEY insight from reviewer Snn46 --
$\hat{\kappa} \to 2/R$ holds \emph{at the isometric fixed point} where the
decoder's pullback metric equals the true metric. The factor $\sqrt{3/G}$ cancels
when the decoder is isometrically calibrated (scale factor $\alpha = \sqrt{3/G}$).

\textbf{Resolution}: Both formulas are correct for different regimes:
\begin{itemize}
  \item Uncalibrated decoder (baseline or early training): $\hat{\kappa} \to (2/R)\sqrt{3/G}$
  \item Isometrically calibrated decoder (at fixed point): $\hat{\kappa} \to 2/R$
\end{itemize}
This makes $\hat{\kappa}$ a DUAL metric: curvature detector AND isometry quality
measure. The experimental target $\hat{\kappa}^* = 2.0$ (used in our Table 1)
is the \emph{asymptotic isometric target}, not the uncalibrated prediction.

\textbf{Action}: Remove the current Eq.~(16) (which incorrectly claims it is
``scale-invariant'' for all decoders) and replace with a Proposition stating:
\begin{enumerate}
  \item For uncalibrated decoders: $\hat{\kappa} \to (2/R)\sqrt{3/G}$
  \item At the isometric fixed point: $\hat{\kappa} \to 2/R = 2\sqrt{K}$
  \item Corollary: $\hat{\kappa} = 2\sqrt{K}$ is achievable \emph{if and only
    if} the decoder is isometrically calibrated.
\end{enumerate}
This transforms the experimental comparison: SCR-VAE's $1.442$ vs. baseline's
$2.849$ (for sphere, true target $2.0$) shows the SCR-VAE is closer to the
isometric target despite the baseline overshooting due to Euclidean KNN
spurious curvature effects.
\end{correct}

\subsection{W3: Torus Wrong Manifold (Critical)}

\begin{accept}[Accepted -- confirmed critical bug]
Op46MX correctly confirms the standard torus has
$K(\phi) = \cos\phi/[r(R+r\cos\phi)]$ ranging from $+1/3$ to $-1$ for $R=2, r=1$.
The evaluation metric uses the flat formula.

\textbf{Additional quantification by Op46MX}: At the outer equator, the relative
error in the $\theta$-direction metric is $(R+r)^2/R^2 - 1 = 125\%$ for $R=2, r=1$.
This completely invalidates the torus isometry MAE comparison.

\textbf{Fix}: Clifford torus in $\mathbb{R}^4$ as specified. Priority: Critical.
\end{accept}

\subsection{W4: Global Convergence (Major)}

\begin{partial}[Accepted as limitation]
The ``basin of attraction'' question is inherent to non-convex optimization.
A formal condition of the form ``$K_{\max} r_n^{(0)} < c$ implies Euclidean
KNN $\approx$ Riemannian KNN'' would be valuable. We note this as Open
Problem~2.
\end{partial}

\subsection{W5: Edge KL Not in Theoretical Loss (Moderate)}

\begin{accept}[Accepted -- will add to Equation (3)]
Op46MX correctly identifies that the code computes four loss terms while
Equation~(3) in the paper has only three. The edge KL term
$\beta_e \mathcal{L}_{\mathrm{edge\_KL}}$ (with $\beta_e = 10^{-3}$) acts as
a variational regularizer on the edge encoder $h_\psi$.

\textbf{Effect on theory}: Under $\beta_e \ll 1$, the edge KL is perturbative
and does not affect the fixed-point characterization (Theorem~1) at leading
order. However, it should be acknowledged in Equation~(3) for honesty.

\textbf{Action}: Add the edge KL term to Equation~(3) with a note that
$\beta_e = 10^{-3}$ makes it perturbative.
\end{accept}

\section{Discrepancy D3: Sphere Sampling Not Uniform}

\begin{accept}[Accepted -- minor but will be fixed]
Op46MX correctly notes that using $\theta \sim \mathrm{Uniform}[0, \pi]$ gives
non-uniform density on $S^2$ (true uniform requires $\theta = \arccos(1-2U)$).
This technically violates Assumption~A3 (density bounded below by $\rho_{\min}$).

In practice, the Euclidean projection still gives $\rho_{\min} > 0$ everywhere
(the non-uniformity is not severe for our purposes), but we will fix this
for theoretical consistency.

\textbf{Action}: Update \texttt{sphere()} in \texttt{data/synthetic.py} to
use $\theta = \arccos(1-2U)$ for uniform sampling.
\end{accept}

\section{Bridging Gaps -- Evaluation of Op46MX's Proofs}

\subsection{Lemma 1 (Sobolev Interpolation)}

\textbf{Mathematical evaluation}: \textbf{Correct and clean}. The proof correctly
uses the Taylor expansion of $f_\theta$ and $f_{\mathrm{true}}$ at adjacent points
to bound the Jacobian difference via the reconstruction error. The corollary is
an immediate consequence. This is the cleanest version of this argument among
all four reviewers.

\textbf{Decision}: Incorporate as Lemma and Corollary after Assumption A4'.

\subsection{Corrected Scaling Law for Curvature Proxy (Proposition 2)}

\textbf{Mathematical evaluation}: The computation is correct. The key formula:
\[
  \hat{\kappa} = \frac{\|Ac\|}{\|A\ell_{ij}\|\|A\ell_{ik}\|}
  \approx \frac{\sigma\sqrt{G}\|c\|}{(\sigma\sqrt{G})^2 \|\ell_{ij}\|\|\ell_{ik}\|}
  = \frac{\hat{\kappa}^{d+1}}{\sigma\sqrt{G}}
\]
correctly shows the $1/(\sigma\sqrt{G})$ factor that the paper's Eq.~(16) omits.
For $\sigma = 1/\sqrt{3}$: $\hat{\kappa} = \hat{\kappa}^{d+1} / (\sqrt{G/3})$.

\textbf{Nuance}: This formula holds for \emph{uncalibrated} decoders. At the
isometric fixed point (from Snn46's analysis), the denominator scales differently
because the log maps are rescaled to match the true metric, canceling the
$1/\sqrt{G/3}$ factor.

\textbf{Decision}: Include both formulas (uncalibrated and calibrated) in the paper.

\subsection{Isometry Lower Bound for Torus (Section 5.3)}

\textbf{Mathematical evaluation}: The lower bound $\max|d_R^* - d_M| \geq \pi r$
follows from the Borsuk-Ulam-type argument: any continuous $g: T^2 \to \mathbb{R}^2$
must identify antipodal points on the shorter generating circle (length $2\pi r$).
The proof sketch is valid.

\textbf{Note}: The bound $\pi r = \pi$ (for $r=1$) is the \emph{worst-case}
error; the average-case MAE of $\approx 1.35$ is indeed below this bound because
random pair sampling avoids the worst-case antipodal pairs.

\textbf{Decision}: Add this as a Proposition to the paper.

\subsection{Multi-scale Iteration for Torus}

\begin{partial}[Interesting but deferred]
Op46MX proposes a coarse-to-fine KNN strategy (large $k$ initially, decreasing
over iterations). This is a practical improvement but requires algorithmic
changes. We defer to future work.
\end{partial}

\section{Op46MX Follow-Up: Correction on Curvature Proxy Target}

Op46MX disputes the dual-regime framing adopted from Snn46 and identifies a
physical error: the "isometric calibration" condition $\alpha = \sqrt{3/G}$
is unreachable under MSE training.

\begin{correct}[Critical correction accepted -- Op46MX is right]
Snn46's algebra is correct: substituting $\alpha = \sqrt{3/G}$ into
$\hat{\kappa} = (2/(R\alpha))\sqrt{3/G}$ gives $\hat{\kappa} = 2/R$. However,
the premise $\alpha = \sqrt{3/G} \approx 0.245$ requires the decoder to output
4$\times$ smaller than the data (since $\|Af_{3D}\| \approx \sqrt{G/3} \approx 4$
for $G=50$). Under MSE training the reconstruction loss $\approx (1-\alpha)^2
\|Af_{3D}\|^2 \approx 9.5$ at this scale -- enormous compared to the observed
$\approx 0.01$. MSE training drives $\alpha \to 1$, not $\alpha \to \sqrt{3/G}$.

\textbf{Correct comparison target:} For $\alpha = 1$ (data-matching decoder,
which is what MSE training produces): $\hat{\kappa}_{\mathrm{theory}} =
(2/R)\sqrt{3/G} = 0.4899$ for $R=1$, $G=50$.

\textbf{Both models overshoot 0.4899} (baseline 2.849, SCR-VAE 1.442) because
finite triangle sizes and non-local graph edges inflate the closure vectors
beyond the $\varepsilon \to 0$ prediction. The improvement from using
Riemannian KNN (local edges) vs. Euclidean KNN (non-local edges) is:
\begin{align*}
  \text{Baseline deviation:} &\quad |2.849 - 0.489| = 2.360 \\
  \text{SCR-VAE deviation:} &\quad |1.442 - 0.489| = 0.953 \\
  \text{Improvement:} &\quad 59.6\% \quad (\text{was incorrectly reported as } 34\%)
\end{align*}
The experimental result is actually \emph{stronger} under the correct comparison.

\textbf{Actions}:
\begin{enumerate}
  \item Update Table~1: change ``True'' column for $S^2$ from $2.000$ to $0.4899$.
  \item Update curvature proxy prediction paragraph to use the reconstruction-regime
    formula $(2/R)\sqrt{3/G}$ as the experimental comparison target.
  \item Add note: $2/R$ is the asymptotic isometric limit (unreachable with MSE
    training); $(2/R)\sqrt{3/G}$ is the reconstruction-dominated regime target.
  \item Report improvement as $59.6\%$ in the paper.
  \item Alternatively (Op46MX's option 2): define a normalized proxy
    $\tilde{\kappa} = \hat{\kappa} \cdot \sqrt{G/3}$ and compare to $2/R = 2.0$.
    Both normalizations are equivalent; we will use the raw proxy with target $0.4899$.
\end{enumerate}
\end{correct}

\begin{partial}[Minor concerns -- partially accepted]
\textbf{``SCR-VAE moving toward isometric target'' framing}: Op46MX correctly
notes this is misleading without tracking $\alpha$ separately. Corrected framing:
``SCR-VAE reduces curvature artifacts from non-local Euclidean KNN edges.''

\textbf{Overshoot explanation}: Op46MX notes that $\hat{\kappa} > 0.489$ could
reflect $\alpha < \sqrt{3/G}$ or finite-triangle effects. We agree both are
possible; we add a note in the paper acknowledging this ambiguity. Measuring
$\alpha$ directly (as the ratio of mean log-map norm to theoretical prediction)
would disambiguate, and we will add this as an additional diagnostic metric.
\end{partial}

\section{Action Items from Op46MX}

\begin{tabular}{lll}
\hline
Item & Status & Priority \\
\hline
Fix torus: Clifford embedding in $\mathbb{R}^4$ & Accept & Critical \\
Resolve curvature proxy formula (dual regime) & Accept + significant & Critical \\
Add edge KL to Equation (3) & Accept & Moderate \\
Fix sphere sampling to uniform & Accept & Minor \\
Sobolev lemma for A4' & Accept & High \\
Empirical verification of A4' & Accept & Moderate \\
Isometry lower bound for torus & Accept & Moderate \\
Multi-scale iteration & Defer & Low \\
\hline
\end{tabular}

\end{document}
